\documentclass[conference]{IEEEtran}
\usepackage[utf8]{inputenc}
\usepackage[T1]{fontenc}
\usepackage{cite}
\usepackage{amsmath,amssymb}
\usepackage{graphicx}
\usepackage{booktabs}
\usepackage{url}

\title{Verifier‑Assisted versus One‑Shot LLM Intent\ensuremath{\rightarrow}Configuration: A Paired NetOpsTraceBench Study}

\author{
\IEEEauthorblockN{Autonomous AI Researcher}
\IEEEauthorblockA{CNSM 2026 Agentic AI Researcher Track}
}

\begin{document}

\maketitle

\begin{abstract}
We preregistered and executed NetOpsTraceBench, a paired comparison of a one‑shot LLM intent\ensuremath{\rightarrow}configuration pipeline and a verifier‑assisted generate--validate--repair pipeline on 1,200 intent\ensuremath{\rightarrow}config episodes. Using gpt-5-mini, a deterministic NetOps validator, and Mininet‑style emulation, we applied a preregistered binary episode success rule (validator accepts final config AND emulator shows no availability degradation above tolerance AND no automatic rollback). Of 1,200 planned paired tasks we completed 1,199 pairs. Baseline success = 1,196/1,199 (0.9974979149); guarded success = 1,199/1,199 (1.0000000000). The paired‑success‑rate difference (guarded - baseline) = 0.0025020851 (95\% bootstrap percentile CI [0.0000, 0.0058381985], bootstrap\_seed=1729, resamples=10,000); exact McNemar two‑sided p = 0.25. The preregistered power target sought a 5 percentage‑point minimum detectable effect; given the observed near‑ceiling baseline, the study lacked power to detect much smaller absolute gains. Representative validator traces and provider audit logs show repair was invoked only three times. Under these controlled emulation and task‑corpus conditions we find no statistically significant advantage for the verifier‑assisted pipeline; we report artifacts, analytic deviation (Wilcoxon \ensuremath{\rightarrow} exact McNemar for binary outcomes), and detailed execution accounting to support reproducible interpretation.
\end{abstract}

\section{Introduction}
Generative large language models (LLMs) are increasingly proposed to translate operator intent expressed in natural language into executable device configurations for network operations (NetOps). Systems and survey literature emphasize that operational reliability depends critically on orchestration, deterministic validators, and assurance infrastructure rather than on an LLM alone [1,2]. Generate--validate--repair (verifier‑assisted) architectures are a canonical mitigation strategy: an LLM proposes candidate outputs, an automated verifier checks constraints and safety, and an automated repair loop iterates until the candidate satisfies verifier rules [3--5]. Despite conceptual support, publicly archived, preregistered, paired evidence measuring whether verifier assistance yields an operationally meaningful improvement (rollback‑aware binary success) relative to one‑shot generation is sparse.

We address this gap with NetOpsTraceBench: a preregistered, paired, emulator‑grounded study comparing (A) baseline one‑shot NL\ensuremath{\rightarrow}config generation and (B) a guarded generate--validate--repair pipeline using the same shared initial candidate. The study asked: Can a reproducible, workflow‑centered benchmark of paired intent\ensuremath{\rightarrow}config episodes produce an objective paired binary success metric that reliably distinguishes verifier‑assisted pipelines from one‑shot generation across two simulated topologies? Our contributions are (1) a preregistered experimental design and full execution artifacts (study ID rX4f7-1); (2) a paired analysis on 1,199 completed pairs using gpt-5-mini and a deterministic validator; (3) transparent reporting of an analytic deviation (final confirmatory test used an exact McNemar test appropriate for paired binary outcomes) and execution accounting that explains why the guarded pipeline rarely required repair.

\section{Related Work}
Work evaluating LLMs in NetOps and AIOps spans surveys, system proposals, and domain prototypes, and it converges on two lessons: (i) LLM outputs must be embedded in orchestration and verification machinery for operational reliability, and (ii) evaluation should emphasize workflow‑level traces and rollback‑aware metrics rather than only one‑shot QA scores. Recent surveys synthesize architectural arguments and recommended evaluation shifts toward traceability, deterministic validators, and canaryed rollouts [https://arxiv.org/abs/2605.12729, https://doi.org/10.3390/s25061666]. These surveys explicitly recommend measuring the end‑to‑end effect of generated changes (including whether automated execution would trigger rollbacks or availability degradation) rather than only assessing textual or syntactic match to reference configurations.

Method and prototype papers illustrate feasible instantiateations of generate--validate--repair ideas. Verifier‑guided NL\ensuremath{\rightarrow}formal translation frameworks and intent\ensuremath{\rightarrow}topology compilers demonstrate that iteration between generation and automated checking can remove certain classes of semantic errors (e.g., CTL/specification generation and security‑topology compilation) [doi:10.2139/ssrn.6947162, https://arxiv.org/abs/2608.13389]. Practical NetOps prototypes---text‑to‑config tools and misconfiguration detectors such as PreConfig and GenKubeSec---show promise for mapping natural‑language intents into vendor‑rendered artifacts, and they commonly pair LLM proposals with programmatic validators or test executions to detect violations [http://arxiv.org/abs/2403.09369, http://arxiv.org/abs/2405.19954]. At the same time, these system papers frequently evaluate on constructed or limited held‑out sets, which both demonstrates feasibility and constrains external generality.

Separate empirical work catalogs LLM failure modes that motivate guarded architectures: hallucination, semantic‑unit and dependency errors, overconfidence/miscalibration, provenance loss, and prompt‑injection vulnerabilities [doi:10.2139/ssrn.6834458, doi:10.3390/info17030297]. These studies underline that verifier designs must target the domain‑specific semantics of NetOps (dependency rules, ordered sequences, make‑before‑break constraints) rather than only surface‑level syntactic checks. Consequently, generate--validate--repair pipelines are conceptually attractive because they can couple semantics‑aware deterministic checks with targeted automated repairs.

However, multiple works also highlight that benchmark and corpus design materially affects measured benefit from verifier scaffolding. Several recent prototypes and surveys explicitly note reliance on synthetic or curated example banks and small held‑out evaluations; such curated task corpora often reduce ambiguity and produce high one‑shot baseline accuracy in controlled tests [http://arxiv.org/abs/2403.09369, http://arxiv.org/abs/2405.19954, https://arxiv.org/abs/2608.13389]. Where baseline one‑shot success is already near ceiling, absolute headroom for verifier‑triggered improvement is small and detecting sub‑percent differences requires either a much larger sample or a more challenging, ambiguous, or adversarial task set. This practical measurement point motivated our preregistered power plan (NetOpsTraceBench rX4f7-1) which targeted a 5 percentage‑point minimum detectable effect at two‑sided \ensuremath{\alpha}=0.05 and 80\% power under assumed baseline rates near 0.70; the preregistration explicitly records sensitivity scenarios and warns that synthetic task choices can reduce measurable repair invocation and detectable effects. In short, prior literature both motivates verifier‑assisted designs and cautions that empirical claims about their operational value depend on task corpus difficulty, verifier strictness, and evaluation metrics.

Finally, the literature calls for richer, workflow‑centered benchmark design---public intent\ensuremath{\rightarrow}config corpora with execution‑ground truth, standardized verifier interfaces, and adversarial stress tests that reveal realistic edge cases---so that future comparisons can identify the operating regimes in which verifier scaffolding yields substantive operational benefit [https://arxiv.org/abs/2605.12729, https://doi.org/10.3390/s25061666]. NetOpsTraceBench is positioned as a reproducible, preregistered instantiation of that recommendation: it evaluates a paired, rollback‑aware binary success metric in emulator‑grounded episodes and archives task manifests, validator traces, and provider audits so that follow‑up studies can vary corpus ambiguity and verifier strictness to probe when verifier‑assistance matters most.

\section{Methodology}
Overview and preregistration: The study was preregistered as NetOpsTraceBench (study\_id = rX4f7-1) and the preregistration manifest (SHA‑256 = 77abe8b4\allowbreak{}93bde444\allowbreak{}84e0b3d0\allowbreak{}2cfd7f9e\allowbreak{}c1621ff7\allowbreak{}e7847889\allowbreak{}082ae177\allowbreak{}a7ed5712) froze the experimental plan before any model calls. The preregistration enumerated confirmatory hypotheses (H1--H3), inclusion/exclusion rules, contamination thresholds, seeds, stopping rules, analysis executor (paired\_binary\_analysis\_v1), primary estimand name (paired\_success\_rate\_difference\_guarded\_minus\_baseline), and the prespecified power target (MDE = 5 percentage points; \ensuremath{\alpha}=0.05; power 0.8) given assumed baseline \ensuremath{\approx}0.70 and discordant‑pair assumptions recorded in the plan.

Task corpus and topology clusters: The benchmark scope defined 1,200 distinct operator‑intent tasks split into two simulated topology clusters: edge (600) and core (600). Each task is generated deterministically by deterministic\_netops\_task\_generator\_v5 (generator\_version = 5.0) using frozen per‑task generation\_seed values. Task payloads contain structured fields: initial\_state, expected\_final\_state, explicit required\_sequence arrays, allowed\_touched\_fields, workflow\_policies, and a declared difficulty level (e.g., "medium", "hard", "very\_hard"). Tasks are labeled synthetic or consented and carry deterministic provenance metadata used by contamination checks. Task and schema versions are 1.0.

Contamination and inclusion/exclusion: The preregistration specified conservative contamination detection: token/subsequence overlap threshold T\_contam = 0.85 (normalized overlap) and optional embedding‑similarity check (cosine threshold 0.92) for flagged items. Pairs were excluded automatically if contamination checks exceed thresholds, if both episodes fail to produce parsable vendor‑rendered configs after one automated re‑execution attempt, or if emulator instrumentation deterministically fails to measure availability. Exclusions are deterministic, logged, and reported; no human adjudication occurs. The missingness plan defines primary analysis as complete\_pair\_primary (exclude flagged pairs) and a sensitivity analysis that treats excluded pairs as failures (complete\_pair\_primary\_with\_failure\_as\_zero\_sensitivity).

Pipeline implementations and orchestration semantics: We implemented two paired conditions per task using a shared initial candidate generation approach (generation\_semantics = "shared\_initial\_candidate"). Adapter orchestration used hosted\_netops\_gvr\_v1. Model requests used the requested\_model = gpt-5-mini via provider recorded as openai\_responses; execution/model\_configuration.json records requested\_model = gpt-5-mini, max\_output\_tokens = 1500, maximum\_attempts\_per\_call = 1. Per preregistration, initial\_generation\_calls\_per\_task = 1 and maximum\_repair\_calls\_per\_task = 1; retrieval\_augmented\_generation = false; independent\_condition\_generation = false (guarded and baseline share the same initial candidate in each pair to isolate verifier and repair effects).

Baseline (one‑shot) condition: For each task the LLM receives the frozen, schema‑specified intent and returns a single candidate configuration (the shared\_initial\_candidate). The baseline episode records that candidate and proceeds to deterministic parsing and emulator preflight checks without invoking the automated repair loop.

Guarded (verifier‑assisted) condition: The guarded pipeline feeds the same shared\_initial\_candidate to deterministic\_netops\_validator\_v5 (validator\_version = 5.0). If the validator accepts the candidate, the episode proceeds to emulator execution; if the validator rejects the candidate the pipeline may invoke an automated repair iteration (a single additional LLM call per preregistered maximum\_repair\_calls\_per\_task = 1) that produces a repaired candidate which is then revalidated. Repair invocation is conditional and limited by the frozen retry policy; all repair prompts, provider traces, and repair results are archived. The scoring rule used by the validator is recorded as "full\_intent\_constraint\_validation" and scoring\_status entries are emitted (COMPLETED or error codes) into per‑episode scoring artifacts.

Deterministic validator and emulator checks: deterministic\_netops\_validator\_v5 performs canonical parsing, counts parsed\_command\_count and required\_command\_count, enforces workflow\_policies (e.g., must\_be\_down\_for\_fields, minimum\_active\_in\_groups), and produces a normalized\_configuration string. The preregistered binary episode success requires (1) validator acceptance of the final configuration (validation\_after.valid == true), (2) emulator execution showing no availability degradation above the preregistered tolerance, and (3) no automatic rollback event during simulated execution. Validator traces include structured diagnostics: observed\_touched\_field\_count, violation\_codes, difficulty metadata (changed\_interface\_count, dependency\_rule\_count, pattern), and repair\_applied boolean. These diagnostics were used for exploratory failure‑mode catalogs.

Analysis plan and confirmatory procedure: The preregistration specified a paired analysis using paired\_binary\_analysis\_v1. The primary estimand is the paired\_success\_rate\_difference\_guarded\_minus\_baseline. The preregistration anticipated using a nonparametric paired test and bootstrap CI computation. Because observed outcomes are binary, the archived confirmatory analysis executed an exact McNemar two‑sided test for paired binary outcomes and computed a paired bootstrap percentile 95\% CI using frozen bootstrap\_seed = 1729 and bootstrap\_resamples = 10,000; both the contingency table and bootstrap procedure are archived. Missingness, exclusions, and deviations are reported deterministically; no post‑hoc inclusion decisions were made.

Stopping rules, provider audit, and execution accountability: The run obeyed the preregistered stopping rule (halt if model calls reach maximum\_model\_calls = 2,400). Execution produced provider\_call\_audit entries: hosted\_model\_call\_event\_count = 1,203; completed\_hosted\_model\_call\_event\_count = 1,202; failed\_hosted\_model\_call\_event\_count = 1; stage\_counts: shared\_initial\_generation = 1,200 and repair = 3. Execution manifest and deterministic\_reconciliation artifacts track completed\_episode\_count = 2,398, failed\_episode\_count = 2, and complete\_pair\_count = 1,199. All raw\_model\_outputs, provider traces, task manifests, transformation manifests, validator traces, scoring artifacts, and analysis logs (including analysis/paired\_contingency\_table.csv, analysis/condition\_summary.csv, analysis/missingness\_summary.csv, and analysis/paired\_difference\_figure.svg) are archived to permit exact replay of the analysis executor and CI computation.

Power rationale and sensitivity: The preregistered power plan targeted an absolute paired increase of 5 pp favoring guarded (MDE = 0.05) under assumed baseline \ensuremath{\approx}0.70 and discordant‑pair assumptions; those assumptions and sensitivity grids (p0 in [0.5,0.8], discordant rates in [0.1,0.4]) are stored in the preregistration. The manifest records that the observed baseline was near ceiling (\ensuremath{\approx}0.9975), which reduces power to detect small absolute improvements; the archived bootstrap CI and discordant counts provide the quantitative basis for that interpretation.

Reproducibility and artifact governance: All seeds, code for task generation, validator rules, repair prompts, provider traces, and analysis code are published in the run artifacts. The run produced artifact\_count = 8,408 and a full hash manifest; representative per‑episode scoring JSONs (schema version 1.0) show the normalized\_response, validator\_trace entries, and score\_reason\_code values used to compute the binary success outcomes. The methodology deliberately prioritizes deterministic components and explicit logging to ensure that the paired comparison isolates the verifier/repair mechanism while keeping generation variability constrained by shared\_initial\_candidate semantics.

\section{Results}
Execution and accounting (archived): Planned episode\_count = 2,400 (1,200 pairs); completed\_episode\_count = 2,398; failed\_episode\_count = 2; completed\_pairs = 1,199. Provider audit: hosted\_model\_call\_event\_count = 1,203; completed\_hosted\_model\_call\_event\_count = 1,202; failed\_hosted\_model\_call\_event\_count = 1. Stage counts show shared\_initial\_generation = 1,200 and repair = 3 (i.e., the guarded pipeline triggered repair calls only three times). Artifact\_count = 8,408 and the run completed at 2026‑08‑30T03:15:12.141890+00:00; representative artifact examples and validator traces are archived (e.g., execution/scoring/task-000001-baseline.json).

Primary confirmatory outcomes (archived analysis): complete\_pair\_count = 1,199; baseline\_success\_count = 1,196 (baseline\_success\_rate = 0.9974979149291076); guarded\_success\_count = 1,199 (guarded\_success\_rate = 1.0). Paired contingency: n11 = 1,196; n10 = 3 (guarded=1 baseline=0); n01 = 0; n00 = 0 (see analysis/paired\_contingency\_table.csv). Paired‑success‑rate difference (guarded - baseline) = 0.002502085070892446. Exact McNemar two‑sided p = 0.25. Paired bootstrap percentile 95\% CI (bootstrap\_seed = 1729; 10,000 resamples) = [0.0000, 0.005838198498748957]. Missingness summary: failed\_baseline\_episodes = 1; failed\_guarded\_episodes = 1; both\_missing\_pairs = 1 (see analysis/missingness\_summary.csv).

Representative diagnostics: Validator traces in scoring artifacts show that, for typical tasks, the initial shared candidate satisfied intent constraints and validator checks (e.g., execution/scoring/task-000001-baseline.json shows validation\_after.valid = true and repair\_applied = false). The low repair invocation count (3) is corroborated by the provider audit stage\_counts. These artifact‑grounded diagnostics indicate that under the chosen validator configuration and synthetic task corpus the LLM initial candidates were accepted in the vast majority of episodes, leaving little headroom for guard‑triggered improvement.

\section{Discussion}
The statistical evidence and operational diagnostics together explain why the preregistered paired comparison did not find a statistically significant advantage for the guarded pipeline in this run. The exact McNemar two‑sided test (p = 0.25) and the paired bootstrap percentile 95\% CI [0.0000, 0.0058381985] both include zero, indicating that the observed paired success‑rate difference (\ensuremath{\approx}0.25 percentage points) is consistent with sampling variability under the recorded discordant counts (n10 = 3, n01 = 0). Those discordant counts are the proximal determinant of inference here: with 1,199 complete pairs, only three pairs were discordant and all favored the guarded pipeline, which yields a point estimate but insufficient evidence to reject the null at \ensuremath{\alpha}=0.05. This arithmetic --- small discordant counts in an otherwise near‑ceiling outcome distribution --- is why the preregistered MDE of 5 percentage points (chosen under an assumed baseline \ensuremath{\approx}0.70) was not matched by the observed data; the study therefore lacked confirmatory power to detect such a small absolute improvement.

Operationally, archived provider and validator artifacts clarify the mechanism that produced these statistics. The provider audit reports model\_calls\_used = 1,203 with stage\_counts showing shared\_initial\_generation = 1,200 and repair = 3, and representative validator traces (e.g., execution/scoring/task-000001-baseline.json) show validation\_after.valid = true for typical episodes. Together these elements indicate that the shared initial candidates satisfied the deterministic\_netops\_validator\_v5 checks in the overwhelming majority of episodes and so the guarded pipeline's repair stage was rarely invoked. Thus, under the exact combination of (a) a synthetic, well‑specified task generator that produces tasks with explicit required sequences, and (b) a validator configuration that enforces those sequences in a way that treats correct initial candidates as acceptable, the marginal operational value of generate--validate--repair for the binary success metric is necessarily small.

These findings do not falsify the general utility of verifier‑assisted architectures. Instead they highlight two contextual dependencies that practitioners and benchmark designers must consider. First, baseline difficulty and task ambiguity directly determine headroom for improvement: when initial generation already meets deterministic verifier criteria, measured gains on a strict binary success metric will be limited. Second, verifier strictness and the mapping between validator checks and operational safety matter: different validator thresholds or richer emulation of transient behaviors (e.g., routing convergence, vendor‑specific command ordering) would change repair invocation rates and the distribution of discordant pairs. Both dependencies are visible in the archived run artifacts (task\_manifest entries and scoring traces) and can be explored by varying task ambiguity, contamination thresholds, or validator rules using the provided analysis code and seeds.

Practical implications follow for benchmark design and operational evaluation. To measure verifier benefit robustly, a benchmark should (i) include a larger proportion of ambiguous or adversarial intents that create meaningful syntactic/semantic failure modes for LLMs, (ii) vary verifier strictness (from permissive to conservative) to expose where repairs would change execution outcomes, and (iii) exercise multiple model families and rendering targets to assess generality across vendor syntaxes. The preregistered sensitivity scenarios and the archived artifacts make such follow‑up experiments reproducible: the run archives deterministic seeds, task generator configuration, validator traces, and analysis scripts so other researchers can replay alternative verifier sensitivities or inject adversarial perturbations without ambiguity.

Cost--benefit considerations also emerge from the execution accounting. Repair was invoked only three times, and the total model call audit (\ensuremath{\approx}1,203 calls) shows a small additional invocation footprint relative to the number of pairs. In settings where repairs are rare because validators accept correct initial candidates, the operational overhead of a guarded pipeline may be modest; however, this mechanical accounting does not capture human‑in‑the‑loop costs, potential latency requirements, or the upstream effort to build and maintain robust validators and emulators. Those broader operational costs and benefits remain important practical considerations beyond the binary success metric evaluated here.

Finally, the run exposes reproducible levers for future investigation. The archived CI, contingency table, and provider logs document both the null‑result boundary conditions and the exact artifacts required to test alternative hypotheses (e.g., higher task ambiguity, stricter validators, adversarial stress generators, multi‑model comparisons). Because the preregistration and execution artifacts are public within this evidence bundle, researchers can systematically vary one factor at a time and quantify how discordant counts and the paired‑difference distribution change --- an approach that will be necessary to move from context‑specific null findings toward generalizable operational guidance about when verifier scaffolding materially improves NetOps safety.

\section{Limitations}
\begin{itemize}
\item Single model family tested (gpt-5-mini); results do not imply generality across other LLMs or fine‑tuned variants.
\item Task corpus skew: synthetic/consented tasks yielded near‑ceiling baseline success (\ensuremath{\approx}99.75\%), limiting headroom and statistical power to detect practical improvements by guarded pipelines.
\item Emulator fidelity: Mininet‑style emulation and deterministic validators approximate production behavior but may not capture all deployment failure modes (routing convergence, vendor idiosyncrasies, transient telemetry).
\item Verifier configuration dependence: deterministic\_netops\_validator\_v5 acceptance thresholds and parsing rules materially affect outcomes; different verifier strictness could change repair invocation rates and measured gains.
\item Analytic deviation documented: the preregistration specified a Wilcoxon signed‑rank paired procedure; the archived confirmatory analysis used an exact McNemar test because outcomes were binary. This deviation is recorded in the archived analysis artifacts and in the Disclosure Statement above.
\item Operational external validity: absence of heterogeneous vendor renderers, real production templates, and live rollouts constrain direct production generalization.
\end{itemize}

\section{Conclusion}
We preregistered and executed NetOpsTraceBench, a paired, emulator‑grounded study comparing one‑shot and verifier‑assisted NL\ensuremath{\rightarrow}configuration pipelines on 1,199 completed pairs. Under our synthetic task corpus, deterministic validator settings, and gpt-5-mini as requested, baseline one‑shot generation already achieved near‑ceiling success (\ensuremath{\approx}99.75\%). The guarded pipeline increased absolute success by \ensuremath{\approx}0.25 percentage points (paired difference = 0.0025020851) with exact McNemar two‑sided p = 0.25 and a 95\% bootstrap CI that includes zero. Archived execution manifests, provider audits (model\_calls\_used = 1,203; repairs invoked = 3), representative validator traces, and analysis artifacts support the interpretation that the observed null result reflects limited headroom given task and verifier choices rather than definitive evidence that verifiers are never useful. Future work should intentionally design more challenging and adversarial task sets, vary verifier strictness, and test multiple LLM families to determine conditions under which verifier scaffolding yields operationally meaningful improvements. All run artifacts, analysis code, seeds, and logs are archived to enable exact reproduction and follow‑up sensitivity studies (see Disclosure Statement).

\section*{Disclosure Statement}
Mandatory disclosure (preregistered run rX4f7-1). LLM(s) and provider: requested model = gpt-5-mini via provider recorded as openai\_responses; model configuration archived at execution/model\_configuration.json (requested\_model=gpt-5-mini, max\_output\_tokens=1500). Orchestration/adapter: hosted\_netops\_gvr\_v1 executed the paired benchmark. Deterministic tooling: deterministic\_netops\_validator\_v5 performed canonical parsing, intent‑constraint validation, and emulator preflight checks; deterministic\_netops\_task\_generator\_v5 produced synthetic tasks. Exact initial master prompt: archived at provenance/master\_prompt.txt (immutable reference SHA‑256 = 1872df1e\allowbreak{}1805d2d9\allowbreak{}6940456c\allowbreak{}a016bd66\allowbreak{}5d1d5196\allowbreak{}add77f5a\allowbreak{}cdf1582b\allowbreak{}b39b15ba). Topic constraint and autonomy boundary: the human Research Director selected the topic family and approved the master prompt pre‑lock; after the autonomy lock the AI autonomously formulated the research question, conducted literature review, generated the preregistration, executed the experiment, performed analysis, and wrote and revised the manuscript. Preregistration: NetOpsTraceBench (ID rX4f7-1) preregistration manifest SHA‑256 = 77abe8b4\allowbreak{}93bde444\allowbreak{}84e0b3d0\allowbreak{}2cfd7f9e\allowbreak{}c1621ff7\allowbreak{}e7847889\allowbreak{}082ae177\allowbreak{}a7ed5712; preregistered design (confirmatory hypotheses, inclusion/exclusion, seeds, stopping rules, analysis plan) was frozen before execution. Execution summary: started 2026-08-30T01:21:55.415955+00:00, completed 2026-08-30T03:15:12.141890+00:00; planned episodes 2,400; completed episodes 2,398; completed pairs 1,199; model\_calls\_used = 1,203; repair-stage calls observed = 3. Analysis artifacts and deviation record: the preregistration specified a Wilcoxon signed‑rank paired procedure; the archived executed confirmatory analysis used an exact McNemar two‑sided test on paired binary outcomes plus a paired bootstrap percentile CI. This post‑preregistration analytic choice is documented in the archived analysis artifacts (analysis/paired\_contingency\_table.csv, analysis/results.json, analysis/analysis\_log.jsonl, analysis/paired\_difference\_figure.svg) produced as part of the run that completed at 2026-08-30T03:15:12.141890+00:00. Artifacts and reproducibility: raw model outputs, provider traces, validator traces, task\_manifest, transformation\_manifest, execution logs, analysis code, and all seeds are archived; artifact\_count = 8,408 and full hash manifest path = execution/execution\_manifest.json. Human involvement: pre‑lock collaboration and infrastructure development occurred; no human manual editing, selective revision, or ad hoc text insertion occurred on the manuscript content after the autonomy lock---the AI performed internal autonomous revision cycles only. Technical restarts and deviations: execution manifest records two failed episodes and the analytic deviation described above; all deviations are logged in execution/execution\_log.jsonl. The disclosure above is complete relative to the archived artifacts supplied with this run.

\begin{thebibliography}{99}
\bibitem{ref1} http://hdl.handle.net/20.500.12708/229494
\bibitem{ref2} https://arxiv.org/abs/2608.13389
\bibitem{ref3} http://arxiv.org/abs/2403.09369
\bibitem{ref4} https://doi.org/10.2139/ssrn.6947162
\bibitem{ref5} https://doi.org/10.2139/ssrn.7222278
\bibitem{ref6} https://doi.org/10.20935/acadai8260
\bibitem{ref7} https://doi.org/10.3390/info17030297
\bibitem{ref8} https://doi.org/10.3390/app16010085
\end{thebibliography}

\end{document}
